% !TEX root = ../../../main.tex
\section{Eigenvalues and eigenspaces}
\label{sec:linear-algebra-i:eigenvalues-and-eigenspaces}

% Sources: IMG_9267--IMG_9269, IMG_9317--IMG_9320;
% October 15 PDF, pp. 3--5.
\begin{definition}
  Let \(L:V\to V\) be linear. A scalar \(\lambda\in\mathbb F\)
  is an eigenvalue of \(L\) if there exists \(v\in V\setminus\{0\}\)
  such that
  \[
    L(v)=\lambda v.
  \]
  Such a nonzero vector \(v\) is an eigenvector for \(\lambda\).
\end{definition}

The identity map \(I\) allows us to subtract scalar multiplication from
\(L\). Since \((\lambda I)(v)=\lambda v\),
\[
  L(v)=\lambda v
  \quad\Longleftrightarrow\quad (L-\lambda I)(v)=0
  \quad\Longleftrightarrow\quad v\in\ker(L-\lambda I).
\]
For a nonzero solution, \(L-\lambda I\) is singular. Thus
\autoref{thm:linear-algebra-i:determinant-invertibility} gives
\begin{equation}
  \lambda\text{ is an eigenvalue of }L
  \quad\Longleftrightarrow\quad \det(L-\lambda I)=0.
  \label{eq:linear-algebra-i:eigenvalue-determinant-criterion}
\end{equation}
% IMG_9268 ends with a nonzero determinant. IMG_9320 and the PDF
% explicitly give zero, consistently with singularity.

\begin{definition}
  The eigenspace associated with \(\lambda\) is
  \[
    E_\lambda(L)=\ker(L-\lambda I)
      =\{v\in V\mid L(v)=\lambda v\}.
  \]
\end{definition}

An eigenspace includes the zero vector, while an eigenvector must be
nonzero. In particular, \(E_0(L)=\ker L\). A nonzero scalar multiple
of an eigenvector is another eigenvector for the same eigenvalue.

\begin{proposition}
  Every eigenspace is a subspace of \(V\). If \(\lambda\) is an
  eigenvalue, then \(\dim E_\lambda(L)>0\).
\end{proposition}

\begin{proof}
  The zero vector belongs to \(E_\lambda(L)\). For
  \(v_1,v_2\in E_\lambda(L)\) and \(a_1,a_2\in\mathbb F\),
  \begin{align*}
    L(a_1v_1+a_2v_2)
      &=a_1L(v_1)+a_2L(v_2)\\
      &=a_1\lambda v_1+a_2\lambda v_2\\
      &=\lambda(a_1v_1+a_2v_2).
  \end{align*}
  Hence \(a_1v_1+a_2v_2\in E_\lambda(L)\). For an eigenvalue,
  the eigenspace contains a nonzero vector, so its dimension is positive.
\end{proof}

\begin{proposition}
  \label{thm:linear-algebra-i:distinct-eigenspaces-intersection}
  If \(\lambda_1,\lambda_2\) are distinct eigenvalues of \(L\), then
  \[
    E_{\lambda_1}(L)\cap E_{\lambda_2}(L)=\{0\}.
  \]
\end{proposition}

\begin{proof}
  If \(v\) belongs to both eigenspaces, then
  \(L(v)=\lambda_1v\) and \(L(v)=\lambda_2v\). Subtracting gives
  \((\lambda_1-\lambda_2)v=0\). Since the scalar is nonzero,
  \(v=0\).
\end{proof}

% !TEX root = ../../hw06.tex
% Source: supplied Homework 6 PDF, page 1, Question 4.
\begin{exercise}
  \label{ex:hw06:sums-of-eigenvectors}
  Suppose, for a matrix \(A\), \(v_1\) is an eigenvector for the eigenvalue
  \(\lambda_1\) and \(v_2\) is an eigenvector for the eigenvalue
  \(\lambda_2\). What are the necessary and sufficient conditions that
  \(v_1+v_2\) be an eigenvector of \(A\) and what is its eigenvalue?
  \solutionlink{hw06:sums-of-eigenvectors}
\end{exercise}


\subsection{Matrices and coordinates}

For \(M\in\mathbb F^{n\times n}\), a nonzero column \(a\in\mathbb F^n\)
is an eigenvector with eigenvalue \(\lambda\) if \(Ma=\lambda a\).
The eigenspace is \(E_\lambda(M)=\ker(M-\lambda I)\), and the
eigenvalue criterion is \(\det(M-\lambda I)=0\).

Let \(E=(e_1,\ldots,e_n)\) be an ordered basis and let \(L(E)=EM\).
For \(v=Ea=\sum_{j=1}^n e_ja^j\),
\begin{align*}
  L(v)&=\sum_{j=1}^n a^jL(e_j)\\
    &=\sum_{j=1}^n a^j\sum_{i=1}^n e_iM^i_j\\
    &=\sum_{i=1}^n e_i\left(\sum_{j=1}^n M^i_ja^j\right)
      =E(Ma).
\end{align*}
On the other hand, \(\lambda v=E(\lambda a)\). Uniqueness of coordinates
therefore gives
\[
  L(v)=\lambda v\quad\Longleftrightarrow\quad Ma=\lambda a.
\]
Thus a nonzero vector is an eigenvector of \(L\) precisely when its
coordinate column is an eigenvector of the representing matrix.

\subsection{A motivating quotient}

% IMG_9267 also asserts that critical points are eigenvectors without a
% symmetry hypothesis. On October 22 Lucas asked to leave that claim flagged.
% Retain the supplied quotient, but do not assert the critical-point claim
% or identify its supremum with an eigenvalue until the hypothesis is clarified.
For a real matrix \(M\) and a nonzero vector \(v\in\R^n\), the quotient
\[
  \frac{v\cdot Mv}{v\cdot v}
    =\frac{\displaystyle\sum_{j,k=1}^n v^jM^j_kv^k}
      {\displaystyle\sum_{j=1}^n(v^j)^2}
\]
compares a quadratic expression with the squared length of \(v\).
The associated optimization problem is to consider
\[
  E(M)=\sup_{v\in\R^n\setminus\{0\}}
    \frac{v\cdot Mv}{v\cdot v}.
\]
